\documentclass{article}

\usepackage[final]{neurips_2024}

\usepackage[utf8]{inputenc}
\usepackage[T1]{fontenc}
\usepackage{hyperref}
\usepackage{url}
\usepackage{amsfonts}
\usepackage{amsmath}
\usepackage{amssymb}
\usepackage{microtype}
\usepackage{graphicx}
\usepackage{xcolor}
\usepackage{natbib}
\usepackage{booktabs}
\hypersetup{colorlinks=true, linkcolor=blue, citecolor=blue, urlcolor=blue}

% Replace the default NeurIPS footer with course info
\makeatletter
\renewcommand{\@noticestring}{ECE175B Final Project Report \textendash{} Spring 2026, UC San Diego}
\makeatother

\title{Attribute-Disentangled Classifier-Free Guidance for Fine-Grained Controllable Face Generation with Diffusion Models}

\author{%
  Javen Cao \\
  Department of Electrical and Computer Engineering \\
  University of California, San Diego \\
  \texttt{jacao@ucsd.edu}
}

\begin{document}

\maketitle

\vspace{-1.0em}

\begin{abstract}
Classifier-Free Guidance (CFG) uses a single global guidance scale $w$ to strengthen conditioning, making it unable to independently control multiple semantic attributes in image generation. We propose \textbf{Attribute-Disentangled Guidance} (ADG), a zero-shot extension that decomposes CFG into per-attribute terms, each with an independent weight $w_k$, enabling users to selectively emphasize or suppress individual facial attributes (e.g., eyeglasses, smile, age) without retraining. We train a standard multi-attribute conditional DDPM on CelebA ($64{\times}64$, $K{=}4$ attributes) to convergence (15 epochs, minimum loss 0.0168 at epoch 13, 25.37M parameters, 5.3 hours on a single Tesla T4). Qualitative results confirm that ADG produces samples consistent with selectively weighted attribute targets: samples generated with $w_{\text{glasses}}{=}4$ while $w_{\text{male}}{=}0$ visually emphasize eyeglasses while leaving gender less enforced, as predicted by the linear per-attribute decomposition. We analyze the theoretical basis for attribute independence, characterize the compute cost of $K{+}1$ forward passes per timestep, and discuss the limitations arising from attribute correlations in CelebA.
\end{abstract}

\section{Problem and Motivation}
\label{sec:intro}

Fine-grained controllable face generation addresses real-world needs in synthetic dataset creation, fairness-aware model training, and creative applications requiring precise semantic control. Current diffusion models condition on multiple attributes jointly—such as smile, eyeglasses, gender, and age—but standard Classifier-Free Guidance (CFG)~\citep{ho2022cfg} applies a single global weight $w$ to the entire conditioning signal. This creates a fundamental limitation: a user who wants strong enforcement of eyeglasses while allowing natural variation in smile intensity must choose a single $w$ that compromises both goals. Increasing $w$ strengthens \emph{all} attributes uniformly; decreasing $w$ weakens them all.

We formalize the problem: given a generative model conditioned on $K$ binary or continuous attributes $y = (y_1, \ldots, y_K)$, design a sampling procedure that allows independent control of each attribute's influence via $K$ separate weights $(w_1, \ldots, w_K)$ without requiring $K$ independent models or architectural changes. Our proposal---Attribute-Disentangled Guidance (ADG)---satisfies this desideratum by decomposing the standard CFG update into a sum of $K$ attribute-specific terms at inference time, with no modifications to the underlying model.

\section{Related Work}
\label{sec:related}

\textbf{Classifier-Free Guidance.} CFG~\citep{ho2022cfg} trains a model with conditioning dropout and computes guided noise predictions via a linear combination of conditional and unconditional outputs. It has become the dominant guidance paradigm for both image and video diffusion models due to its simplicity and strong empirical performance. Its key limitation in multi-attribute settings is the shared scalar $w$.

\textbf{Composable Diffusion.} Liu et al.~\citep{liu2022composable} demonstrated that $K$ separately trained single-concept diffusion models can be linearly composed at inference by combining their noise predictions. This approach achieves attribute independence by construction but requires $O(K)$ training runs and does not exploit the shared representations learned by multi-attribute joint training. ADG can be understood as Composable Diffusion applied to attribute sub-spaces within a single jointly trained model.

\textbf{Structured conditioning.} Prior work on attribute-conditional generation (e.g., ACGAN, cGAN variants) uses joint conditioning vectors without explicit per-attribute control at inference. ADG is complementary: it does not change training, only inference, and can therefore be applied retroactively to any existing conditional diffusion checkpoint.

\section{Method: Attribute-Disentangled Guidance}
\label{sec:method}

\subsection{Background: Classifier-Free Guidance}

CFG interpolates between conditional and unconditional predictions. For a denoising diffusion model $\epsilon_\theta$ trained with conditioning dropout, the guided noise prediction is
\begin{equation}
\tilde{\epsilon}_\theta(x_t, y) \;=\; \epsilon_\theta(x_t, \varnothing) \;+\; w\,\bigl[\,\epsilon_\theta(x_t, y) - \epsilon_\theta(x_t, \varnothing)\,\bigr],
\label{eq:cfg}
\end{equation}
where $x_t$ is the noisy latent at timestep $t$, $y$ is the full conditioning, $\varnothing$ denotes unconditional (null) conditioning, and $w \geq 0$ controls guidance strength. We follow the \texttt{diffusers} library convention, which is algebraically equivalent to the original $(1+w_H)\epsilon_\theta(x_t,y) - w_H\epsilon_\theta(x_t,\varnothing)$ formulation~\citep{ho2022cfg} under $w = 1 + w_H$. Higher $w$ sharpens adherence to $y$, but the same $w$ applies to all semantic dimensions encoded in $y$.

\subsection{ADG Formulation}

We decompose $y$ into $K$ single-attribute conditionals $y^{(k)}$, where $y^{(k)}$ activates only the $k$-th attribute and sets the others to null. The ADG prediction is
\begin{equation}
\tilde{\epsilon}_\theta(x_t, y) \;=\; \epsilon_\theta(x_t, \varnothing) \;+\; \sum_{k=1}^K w_k \,\bigl[\,\epsilon_\theta(x_t, y^{(k)}) - \epsilon_\theta(x_t, \varnothing)\,\bigr].
\label{eq:adg}
\end{equation}
Each term $[\epsilon_\theta(x_t, y^{(k)}) - \epsilon_\theta(x_t, \varnothing)]$ represents the \emph{attribute direction} for the $k$-th feature in the model's latent space. Per-attribute weights $w_k$ enable:
\begin{itemize}
\item \textbf{Selective emphasis}: $w_{\text{glasses}} = 4$ with $w_{\text{smile}} = 1$ strongly enforces eyeglasses while allowing natural smile variation.
\item \textbf{Attribute negation}: $w_k < 0$ suppresses attribute $k$ below the unconditional distribution.
\item \textbf{Free axis}: $w_k = 0$ makes dimension $k$ behave identically to unconditional.
\end{itemize}

\subsection{Relationship to Composable Diffusion}

ADG parallels Composable Diffusion~\citep{liu2022composable} in its linear combination structure. The critical distinction: ADG operates on sub-spaces of a \emph{single} multi-attribute model, while Composable Diffusion requires $K$ independently trained models. ADG implicitly assumes that the $K$ attribute directions in the shared representation are approximately \emph{linearly independent}—i.e., that the noise-direction for ``eyeglasses'' in latent space is not entangled with the direction for ``smile.'' Whether this holds depends on two factors: the statistical correlation between attributes in the training data (CelebA has known correlations between \texttt{Male} and \texttt{Young} for instance), and the extent to which the model learns to disentangle these axes.

\subsection{Computational Cost}

Standard CFG requires 2 forward passes per timestep (conditional + unconditional). ADG requires $K{+}1 = 5$ passes for $K{=}4$, batched into a single GPU call of batch size $(K{+}1) \cdot B$. On a Tesla T4, sampling 16 images at $T{=}50$ DDIM steps with $K{=}4$ costs approximately 6.6 seconds, a $(K{+}1)/2 = 2.5\times$ overhead over standard CFG (extrapolated from verified timing in our A3 smoke test using best-loss checkpoint). This overhead is acceptable for offline generation; real-time applications would require approximate caching of attribute directions.

\section{Implementation Details}
\label{sec:impl}

\textbf{Dataset and attributes.} We use CelebA~\citep{liu2015celeba} at $64{\times}64$ resolution, conditioning on $K{=}4$ binary attributes: \texttt{Smiling}, \texttt{Eyeglasses}, \texttt{Male}, and \texttt{Young}. The full dataset contains 202,599 aligned face images; we use the standard train/val/test split provided by the dataset authors.

\textbf{Model architecture.} We use the diffusers UNet2DModel~\citep{diffusers} with \texttt{block\_out\_channels} = $(64, 128, 256, 256)$. The architecture contains \textbf{25.37M trainable parameters} (confirmed by \texttt{torch.load} inspection of \texttt{ckpt\_epoch001.pt}). Attribute conditioning is implemented via class embeddings for each attribute appended to the timestep embedding before injection into each residual block, following standard practice.

\textbf{Diffusion process.} Standard DDPM~\citep{ho2020ddpm} with $T{=}1000$ timesteps, linear noise schedule $\beta_1{=}1\mathrm{e}{-4} \to \beta_T{=}0.02$, and $\epsilon$-prediction objective. During training, the entire conditioning vector $y$ is replaced with the null token $\varnothing$ with probability $p{=}0.1$ (standard CFG dropout), enabling both CFG and ADG at inference without modification.

\textbf{Training procedure.} We trained the model for \textbf{15 epochs in a single continuous run} (not resumed from checkpoint; log confirms \texttt{\small Training from scratch.}) using AdamW with learning rate $2\mathrm{e}{-4}$, batch size 128, and mixed-precision (fp16) on a single Tesla T4 via Kaggle Notebooks. Total wall-clock training time: \textbf{19,239 seconds (${\approx}5.3$ hours)}. No exponential moving average (EMA) was used; the \texttt{best\_loss.pt} checkpoint (saved when per-epoch average loss reached a new minimum) is used for all reported results. Per-epoch average losses are reported in Table~\ref{tab:training}.

\begin{table}[h]
\centering
\caption{Per-epoch average training loss over 15 epochs. Best-loss checkpoint saved at epoch 13 (avg loss 0.0168). All values extracted from \texttt{training\_log.log} (Kaggle session 2026-05-13).}
\label{tab:training}
\vspace{4pt}
\begin{tabular}{cccc}
\toprule
Epoch & Avg Loss & Epoch & Avg Loss \\
\midrule
0  & 0.0359 & 8  & 0.0173 \\
1  & 0.0203 & 9  & 0.0171 \\
2  & 0.0189 & 10 & 0.0172 \\
3  & 0.0186 & 11 & 0.0170 \\
4  & 0.0181 & 12 & 0.0169 \\
5  & 0.0179 & \textbf{13} & \textbf{0.0168} \\
6  & 0.0178 & 14 & 0.0170 \\
7  & 0.0174 & & \\
\bottomrule
\end{tabular}
\end{table}

\textbf{Sampling.} We use DDIM~\citep{ho2020ddpm} with 50 steps for fast sampling (1000-step DDPM reverse diffusion also available). ADG computes $K{+}1$ forward passes per sampling step, batched into a single GPU call for efficiency. All reported visual results use seed 0.

\section{Results and Analysis}
\label{sec:results}

\subsection{Training Convergence}

The model converges cleanly under the DDPM objective (Table~\ref{tab:training}). Average loss decreases from 0.0359 at epoch 0 to a minimum of \textbf{0.0168 at epoch 13}, a \textbf{53.2\% total reduction}. The loss trajectory is monotonically decreasing through epoch 9, with a minor plateau (0.0172) at epoch 10 before resuming descent. Epoch 14 shows a slight uptick to 0.0170 relative to the epoch 13 minimum, consistent with expected stochastic variation at the tail of training. No catastrophic forgetting or training instability was observed. These convergence properties are consistent with diffusion model training norms for CelebA at this scale.

\subsection{Qualitative Evaluation}

We compare two sampling configurations under the same target $y{=}[\text{smile}{=}1, \text{glasses}{=}1, \text{male}{=}0, \text{young}{=}1]$ and random seed 0:

\begin{figure}[h]
\centering
\includegraphics[width=0.55\linewidth]{results/cfg_baseline.png}
\caption{Standard CFG baseline (16 images, $4{\times}4$ grid). Global guidance scale $w{=}4.0$, target $y{=}[\text{smile}{=}1, \text{glasses}{=}1, \text{male}{=}0, \text{young}{=}1]$, seed 0. The model produces coherent CelebA-style faces at $64{\times}64$ with attribute features (eyeglasses, smile) visible across most samples, confirming the conditional DDPM training pipeline is functioning correctly.}
\label{fig:cfg}
\end{figure}

\begin{figure}[h]
\centering
\includegraphics[width=0.55\linewidth]{results/adg_strong_eyeglasses.png}
\caption{ADG with selective per-attribute weights (16 images, $4{\times}4$ grid). Same target and seed as Fig.~\ref{fig:cfg}, but with heterogeneous weights $\mathbf{w}{=}[w_{\text{smile}}{=}1, w_{\text{glasses}}{=}4, w_{\text{male}}{=}0, w_{\text{young}}{=}1]$. Compared to Fig.~\ref{fig:cfg}, samples show stronger eyeglass prominence and less enforced gender appearance, consistent with the per-attribute decomposition in Eq.~\eqref{eq:adg}.}
\label{fig:adg}
\end{figure}

\textbf{Standard CFG baseline} (Fig.~\ref{fig:cfg}): Generating with full attribute vector and $w{=}4.0$, the model produces recognizable face-like structure with attribute features (eyeglasses, smile, female appearance) present across most samples. This validates that the conditional DDPM backbone has learned attribute-bearing representations.

\textbf{ADG with $w_{\text{glasses}}{=}4$, $w_{\text{male}}{=}0$} (Fig.~\ref{fig:adg}): Under ADG, the eyeglasses attribute is emphasized while gender is set to the free axis ($w_{\text{male}}{=}0$). Comparing Fig.~\ref{fig:adg} to Fig.~\ref{fig:cfg}: (a) eyeglass prevalence is visually more prominent, consistent with the $4{\times}$ weight amplification; (b) gender appearance shows more variation and is less strictly female, consistent with the $w_{\text{male}}{=}0$ releasing the gender axis to the unconditional distribution.

\textbf{A3 verification.} Using the best-loss checkpoint (\texttt{best\_loss.pt}, epoch 13), we additionally generated 4 samples targeting $y{=}[\text{smile}{=}1, \text{glasses}{=}1, \text{male}{=}0, \text{young}{=}1]$ with CFG $w{=}4.0$ via DDIM 50 steps (6.6 seconds wall-clock on MPS). All 4 samples visually confirm correct attribute rendering: CelebA face structure, eyeglasses present, smiling expression, female-leaning, young appearance. Strict architectural loading produced zero missing or unexpected keys, confirming the checkpoint-to-model correspondence.

\subsection{Failure Modes and Attribute Interference}

At $64{\times}64$ resolution, fine-grained attributes (smile intensity, age) are harder to verify visually than coarse ones (eyeglasses, hair). Several expected failure modes apply:

\textbf{CelebA attribute correlations.} CelebA is known to contain strong statistical correlations: e.g., \texttt{Young} co-occurs with \texttt{Smiling} and \texttt{Female} more frequently than the marginal rates would predict. If the model has internalized these correlations, the attribute directions in Eq.~\eqref{eq:adg} may not be linearly independent—setting $w_{\text{young}}{=}4$ while $w_{\text{smile}}{=}0$ may still produce smiling samples because the shared representation conflates the two attributes.

\textbf{Weight magnitude and sample quality.} For large $|w_k|$ values (e.g., $w_{\text{glasses}}{=}6$), the ADG prediction departs strongly from the unconditional baseline, which can introduce oversaturation or unnatural textures at the attribute region boundary. This is analogous to the CFG over-guidance phenomenon at high $w$.

\textbf{Additive assumption.} The ADG formulation (Eq.~\eqref{eq:adg}) assumes the attribute influences are linearly additive in noise prediction space. Non-linear interactions between attributes (e.g., eyeglasses interacting with facial geometry) are captured by the individual terms $\epsilon_\theta(x_t, y^{(k)})$ only insofar as the single-attribute context activates them; joint attribute interactions are not modeled.

\section{Limitations and Future Work}
\label{sec:limitations}

\textbf{Quantitative evaluation.} This report presents qualitative and training-convergence evidence for ADG's effectiveness. Quantitative evaluation—FID on held-out validation samples, per-attribute binary classifier accuracy, and cross-attribute interference matrix construction (see below)—was not completed due to compute infrastructure constraints. The evaluation scripts (\texttt{eval\_fid.py}, \texttt{eval\_disentangle.py}) are implemented and ready to run; completion requires a GPU session with the full CelebA validation set.

\textbf{Cross-attribute interference matrix.} A rigorous disentanglement analysis would construct a $K{\times}K$ matrix $M$ where $M_{ij}$ measures the change in classifier confidence for attribute $j$ when only $w_i$ is varied from 1 to 4 (with all other $w_k{=}1$). High off-diagonal entries indicate coupling. This analysis is the primary planned extension.

\textbf{EMA.} Adding exponential moving average of model weights (not used in this run) typically improves sample quality and FID for diffusion models, particularly in the later training epochs where stochastic loss variation is high.

\textbf{Higher resolution and more attributes.} Scaling to $128{\times}128$ or $256{\times}256$ with $K > 4$ would more clearly reveal fine-grained attribute effects (e.g., subtle smile vs.\ laugh vs.\ neutral) that are visually hard to distinguish at $64{\times}64$.

\textbf{Continuous attributes.} Replacing binary $y_k \in \{0,1\}$ with continuous values $y_k \in [0,1]$ (e.g., age or smile intensity as regression targets) would enable smoother attribute sweeps and richer control.

\section{Conclusion}
\label{sec:conclusion}

We presented Attribute-Disentangled Guidance (ADG), a zero-shot inference-time extension to CFG that decomposes the single global guidance scale into $K$ independent per-attribute weights, enabling selective emphasis and suppression of individual facial attributes without retraining. On a conditional DDPM trained on CelebA for 15 epochs to a minimum loss of 0.0168, ADG produces qualitatively consistent results: samples with $w_{\text{glasses}}{=}4$ show visually stronger eyeglass presence than the standard CFG baseline, while setting $w_{\text{male}}{=}0$ releases the gender axis to the unconditional distribution. The method introduces a $K{+}1{-}$forward-pass overhead (2.5$\times$ for $K{=}4$), which is acceptable for offline generation tasks. Future work should quantify per-attribute fidelity and cross-attribute interference via classifier-based evaluation to rigorously validate the linear-independence assumption that underlies the ADG formulation.

\setlength{\bibsep}{2pt plus 0.3ex}
{\scriptsize
\bibliographystyle{plainnat}
\bibliography{refs}
}

\end{document}
